%==============================================================================
%  14-data-technology.tex — Data & Technology Foundations
%==============================================================================
\strategyPart{14 \textbar\ Data \& Technology}{Data and Technology Foundations}

AI value is capped by data readiness and architecture. Poor data quality costs
the average organization \textbf{\$12.9~million per year} [Gartner 2021] — the
most direct quantification of the ``data trust deficit.''

\section{Data readiness}

\begin{center}
  \renewcommand{\arraystretch}{1.45}
  \rowcolors{2}{inSnow}{white}
  \fitwidth{\begin{tabular}{>{\raggedright\arraybackslash}p{3.0cm}>{\raggedright\arraybackslash}p{3.6cm}>{\raggedright\arraybackslash}p{3.6cm}>{\raggedright\arraybackslash}p{3.0cm}}
    \rowcolor{inNavy}
    \hcell{Dimension} & \hcell{Low maturity} & \hcell{High maturity} & \hcell{KPI} \\
    \textbf{Quality} & Inconsistent, incomplete & Continuously monitored & Error rate, completeness \\
    \textbf{Access} & Siloed & Governed data products & Access latency \\
    \textbf{Ownership} & Unclear & Enterprise-assigned & Owner per dataset \\
    \textbf{Governance} & Informal & Automated lineage \& controls & Audit coverage \\
  \end{tabular}}
\end{center}

\section{AI architecture layers}

An enterprise agent architecture can be described as nine layers, from
business goals to monitoring [Anthropic; Microsoft agent framework; MCP
documentation]:

\begin{enumerate}
  \item \textbf{User and business-goal layer} — employee and customer
        requests, strategic goals, events and alerts.
  \item \textbf{Interaction layer} — chat, voice, dashboards, business
        applications, collaboration platforms.
  \item \textbf{Orchestration layer} — selects the model and agent, resolves
        required tools, parallelism, approvals, stop and escalation
        conditions.
  \item \textbf{Agent layer} — specialized agents (planning, research, data,
        coding, compliance, service, finance, quality).
  \item \textbf{Model layer} — foundation models, small and domain models,
        predictive, embedding, reranking, vision and speech models.
  \item \textbf{Knowledge and data layer} — databases, warehouses, documents,
        knowledge bases, APIs, ERP/CRM/BI data, product analytics.
  \item \textbf{Tool and action layer} — search, queries, report generation,
        CRM/ERP/ticketing systems, email, calendar, payment and
        product-management platforms.
  \item \textbf{Identity and security layer} — every agent has an identity,
        role, permissions, approved tools, data boundaries, transaction
        limits, audit logs and an emergency stop.
  \item \textbf{Monitoring and evaluation layer} — accuracy, cost, latency,
        tool usage, failed actions, human overrides, policy violations,
        security incidents, satisfaction and business outcomes.
\end{enumerate}

Microsoft's agent framework emphasizes \textbf{human-in-the-loop workflows}:
the agent pauses, sends an approval request and waits before continuing
[Microsoft agent framework]. The OWASP Top 10 for LLM applications flags
\textbf{excessive agency} — granting agents more permissions than a task
requires — as a first-order security risk [OWASP].

\section{Model strategy}

Compare commercial foundation models, open-source models, small language
models, domain models and classical ML across \textbf{accuracy, cost, latency,
privacy, reliability, explainability and vendor dependence}. No single option
wins; the portfolio matters.

\section{Multi-agent integration}

For agentic systems, evaluate: agent registry and capability discovery, task
routing, message formats, shared context, \textbf{agent identity and
inter-agent authentication}, permission boundaries, result validation and
human escalation [MCP, LangGraph, AutoGen, CrewAI documentation].
